\section{从预测损失到组合净收益}
\label{ch:lower-ml-research}

\begin{itemize}
  % Generated from curriculum/manifest.json; do not edit by hand.
\item 直接先修：下册《验证与监控：嵌套时序、校准、解释与漂移》。

\end{itemize}

同一个分数可以对应不同仓位，因而对应不同收益。先指定损失对应的预测对象，再用成本和风险决定是否交易，比直接最大化一条回测曲线更容易解释。


\MFQTransition{预测目标与决策目标}

% Generated by tools/build_figures.py; edit figures/figure-specs.json instead.
\MFQFigure{tex/figures/generated/lower-ml-decision-pnl.pdf}{模型分数经过阈值、仓位上限和成本映射后才形成净收益；预测误差较小并不保证净收益为正。}{lower-ml-decision-pnl}


MSE 衡量点预测误差；pairwise ranking loss 惩罚收益次序与分数次序冲突；收益加权目标直接奖励 $s_ir_i$。三个目标回答不同问题：
\[
L_{\rm rank}=\frac1{|P|}\sum_{(i,j)\in P}
\log\!\left(1+e^{-\operatorname{sign}(r_i-r_j)(s_i-s_j)}\right),
\qquad L_R=-\frac1n\sum_i s_ir_i.
\]
横截面 Alpha 常关心排序，却最终承担成本后的组合收益，预测误差、排序表现与决策收益因此分别回答不同的问题。收益加权损失若不限制杠杆会鼓励无界分数，不能替代组合约束。

\MFQTransition{例：概率与交易阈值}

假设上涨时每单位盈利 $a>0$，下跌时亏损 $b>0$，往返成本为 $c\geq0$，预测上涨概率为 $p$。买入的条件期望净收益为 $pa-(1-p)b-c$，所以仅在
\[
 p>\frac{b+c}{a+b}
\]
时为正。$a=b=1\%$、$c=0.2\%$ 时阈值是 0.6，准确率超过一半仍不够。此推导假定预测概率校准、盈亏幅度固定且成本正确；改变幅度分布后要重新求条件期望。

若连续仓位 $w$ 的收益目标为 $\mu w-\gamma\sigma^2w^2/2-c|w|$，从空仓出发，最优解为 $\operatorname{soft}(\mu,c)/(\gamma\sigma^2)$，再按仓位上限截取。成本使 $|\mu|\leq c$ 的信号不交易。notebook 会画出概率阈值和无交易区，并逐期更新部分成交后的实际仓位。

\MFQTransition{从预测函数到组合策略的六个映射}

这些变量之间的关系可以写为
\[
x_t\xrightarrow{f_\theta}s_t
\xrightarrow{g}\tilde w_t
\xrightarrow{\Pi_\mathcal C}w_t^*
\xrightarrow{\mathrm{orders}}q_t
\xrightarrow{\mathrm{fills}}w_t
\xrightarrow{r_{t+1}}R_{t+1}^{\rm net}.
\]
$s_t$ 是分数，$g$ 可做排序、z-score 或概率到预期收益映射；$\Pi_\mathcal C$ 把目标投影到杠杆、行业、风险和个股约束；订单与成交决定实际仓位；最后才得到净收益。每一步都可能吞掉模型优势。

最简单的横截面方案是按分位数组合。更连续的方案令原始权重与去均值分数成比例：
\[
\tilde w_{i,t}=\frac{s_{i,t}-\bar s_t}
{\sum_j|s_{j,t}-\bar s_t|},
\]
再做行业中性和风险缩放。分母很小时必须拒绝或沿用旧仓位，否则微小分数噪声会被放大成满杠杆。

\MFQTransition{信号滞后、持有期与组合重叠}

预测 $t+1$ 收益不代表能以 $t$ 收盘价成交。必须明确特征截止、模型推理、订单提交和成交价格。若公告在收盘后发布，最早可交易时点通常是下一交易日。日历、停牌和涨跌停会改变持有期。

$h$ 期标签对应的策略若每天建仓，会同时持有 $h$ 个 vintage。组合收益是各 vintage 的叠加，换手也由实际总仓位变化计算；不能把每个标签样本当作独立完整交易再求平均。重叠持有还使统计标准误需要处理序列相关。

\MFQTransition{风险约束与模型不确定性}

目标权重常通过优化得到：
\[
\max_w\quad \hat\mu_t^{\mathsf T}w
-\frac\gamma2w^{\mathsf T}\Sigma_tw
-\kappa\|w-w_{t-1}\|_1,
\]
约束包括净/毛敞口、行业和风格中性、单资产上限及可交易性。机器学习只提供 $\hat\mu_t$ 或排序；协方差、成本和约束来自其他模型。

预测不确定性可以进入 shrinkage 或仓位缩放。例如令
\[
\mu_{i,t}^{\rm adj}=
\frac{\hat\mu_{i,t}}{1+\lambda\hat\sigma_{i,t}},
\]
但 $\hat\sigma$ 必须经过样本外校准。简单把神经网络 softmax 当置信度会产生过度集中。

\MFQTransition{分数、目标仓位与部分成交}

每期做多最高一只、做空最低一只，毛敞口为 2。第一期分数 $[0.9,0.2,-0.8]$ 产生目标 $[1,0,-1]$。第二期目标为 $[-1,1,0]$；从旧仓位到目标的订单增量为 $[-2,1,1]$。成交比例 $[0.5,1,0]$ 后，实际仓位为 $[0,1,-1]$：未成交的第三只仍保留旧空仓，而不是强制归零。

两期毛收益分别为 $0.03$ 与 $0.02$。实际成交换手为 $2+2=4$，单位成本 $0.001$，所以
\[
R^{\mathrm{net}}=0.05-4(0.001)=0.046.
\]
这个计算的关键是逐期更新“旧仓位加成交增量”。未成交部分不会自动消失，目标仓位也不能直接替代实际仓位。配套实验显示每期分数、目标与成交后的持仓，便于改变成交比例后重算收益。这里的比例仍是外部给定的简化假设，没有进一步模拟队列、借券或市场冲击。

\MFQTransition{成本、冲击与容量}

线性成本近似
\[
C_t^{\rm lin}=\sum_i c_{i,t}|\Delta w_{i,t}|
\]
适合佣金和半点差。市场冲击常随交易量占比非线性增长，例如平方根近似
\[
C_{i,t}^{\rm impact}\approx
\sigma_{i,t}|\Delta w_{i,t}|
\eta\sqrt{\frac{Q_{i,t}}{V_{i,t}}}.
\]
它是经验模型，参数需按市场和订单类型校准。容量分析应逐步放大资本，重算订单量、冲击和借券约束，观察净收益在何处归零，而不是把历史百分比收益线性乘资本。

成本模型也必须只用当时可得的流动性。用当天收盘后才完整知道的成交量估算当日开盘订单冲击，是未来信息。

\MFQTransition{模型集成与简单基线}

树、神经网络、文本和图模型可做简单平均、基于验证表现加权或 stacking。元模型只能读取 base model 的 out-of-fold 预测；用同一训练样本的拟合值训练 stacking 会严重乐观。高度相关模型的平均并不会带来多少分散，应报告预测相关和边际贡献。

比较集成模型时，可以保留一条预先选定的树模型或线性模型作为简单基线。复杂模型在某些窗口表现较差，是理解其适用范围的材料；根据最终评价结果反复更换基线，会使比较本身参与模型选择。

\MFQTransition{历史评价与向前观察}

历史回测考察给定数据与成交假设下的表现，向前观察则考察新数据到来时信号怎样变化。在不实际交易的观察中，成交与市场影响仍需假设。因此，不同评价方式各自提供不同信息，不能简单按一种强弱次序替代。

最终测试结果用于修改模型后，这段数据便参与了开发。之后的表现可以解释开发过程，却需要新的未来区间来评价调整后的预测能力。

\begin{diagnostic}
若收益只来自极少数日期、单一行业或无法成交的小盘股，模型指标再漂亮也不是稳健 Alpha。逐笔归因、暴露分解和流动性过滤应在宣称改进之前完成。
\end{diagnostic}

\MFQTransition{文本信息的可得时间}

文本的首次发布与所使用的修订版本都应早于决策时点。例如，1 月 3 日作出预测时，1 月 5 日才发布的修订内容尚不可用，即使该文本最初发布得更早也一样。

预训练数据也可能包含评价期之后的信息。分析历史文本时，除了文本版本，还需考虑模型训练截止与可能的记忆影响。若只能获得当前修订版本，实验可以说明一种分析方法，但难以据此重建历史可用的信息。

\section{本路线的综合项目}

选择一个明确的预测问题，比较一种较复杂的模型与简单基线，并在未参与选择的时间区间评价。将预测转为仓位后，进一步考察部分成交、成本和分布变化怎样影响结果。综合项目应解释预测改善与净收益改善之间的联系。

\section{综合练习}
\begin{enumerate}[leftmargin=*]
  \item \textbf{口述概念：}purge 与 embargo 分别阻断哪条信息路径？
  \item \textbf{口述概念：}内层验证与外层测试为何不能互换？
  \item \textbf{笔试推导：}验证事先确定索引满足两条严格不等式。
  \item \textbf{笔试推导：}写出 Brier 分数并解释校准改善为何不等于收益改善。
  \item \textbf{数值编程：}实现两折时间 cross-fitting，并把预测写回原索引。
  \item \textbf{数值编程：}改变 top-$k$，画出重要性 Jaccard 敏感性。
  \item \textbf{研究判断：}审计在最终测试集上 early stopping 后仍报告“样本外”的研究材料。
  \item \textbf{研究判断：}分别举例说明输入、标签、条件关系和成交变化怎样影响预测或收益，并讨论重新校准与重新估计各能处理什么问题。
\end{enumerate}

\subsection*{基础衔接：独立计算}
真实成功率 $q=0.3$ 时，恒定预测 0.5 比诚实预测多多少期望 Brier 损失？

\MFQLead{Alpha 决策、交易摩擦与研究包题组}
\begin{enumerate}[leftmargin=*]
  \item \textbf{口述概念：}MSE、排序损失与收益加权损失分别服务什么决策？
  \item \textbf{口述概念：}为什么未成交不能把旧仓位直接归零？
  \item \textbf{笔试推导：}重建两期目标、订单增量、实际仓位与 0.046 净收益。
  \item \textbf{笔试推导：}解释无约束收益加权目标为何可能鼓励无界分数。
  \item \textbf{数值编程：}加入仓位上限和借券失败，重算 turnover 与成本。
  \item \textbf{数值编程：}对文本首次发布、抓取与修订时间注入三个泄漏负例。
  \item \textbf{研究判断：}审计一个预测优秀但成本后为负、仍宣称可交易的错误研究包。
  \item \textbf{开放 Capstone：}交付模型、验证、文本协议、组合、摩擦、漂移与限制的一键研究包。
\end{enumerate}

